\section{PMF y el circuito emocional: por qué importa «qué placer»}

La sección anterior trató la presión por sobrevivir; esta sección explica por qué el «qué placer» del título es una señal de producto. No es una simple expresión emocional, sino una de las señales percibidas del PMF. Los usuarios empiezan a asombrarse, el producto empieza a difundirse, la presión de efectivo se alivia, el ánimo del fundador se enciende y el equipo sigue iterando. Este circuito es especialmente visible en las aplicaciones de IA, porque el Magic Moment suele ser muy intenso.

\begin{figure}[H]
\centering
\includegraphics[width=0.96\textwidth]{pmf-emotion-loop.png}
\caption{PMF y el circuito emocional: el reconocimiento de los usuarios, el auge del producto y la satisfacción del fundador se amplifican entre sí. Diagrama conceptual propio, elaborado a partir de la conversación de 01:24:50--01:28:00.}
\end{figure}

\begin{knowledgebox}{Lectura de la figura: la emoción es una señal, pero no lo es todo}
La satisfacción del fundador proviene de la retroalimentación de los usuarios y del alivio de la presión por sobrevivir. Es una señal de que el PMF podría estar apareciendo, pero todavía hay que verificar la retención, el pago, la recompra y la profundidad en el flujo de trabajo.
\end{knowledgebox}

\subsection{Del «qué placer» al negocio}

La sección anterior trató el «placer» como la percepción del PMF; esta sección plantea una pregunta más exigente: cómo se convierte en negocio. En sus primeras etapas, las aplicaciones de IA tienden a difundirse por el asombro que causan, pero un negocio real depende de la retención, el pago y el flujo de trabajo. Un solo «qué placer» indica únicamente que el usuario está dispuesto a probar; el uso continuo indica que el producto entró en el contexto laboral; la disposición a pagar indica que su valor tiene precio.

\begin{figure}[H]
\centering
\includegraphics[width=0.96\textwidth]{magic-moment-to-business.png}
\caption{Del «qué placer» al negocio: el pico emocional debe convertirse en retención, pago y flujo de trabajo. Diagrama conceptual propio, elaborado a partir de la conversación de 01:24:50--01:28:00.}
\end{figure}

\begin{importantbox}{Las tres etapas de validación de una aplicación de IA}
La primera etapa es el asombro: el usuario piensa «de verdad puede hacerlo». La segunda es la retención: el usuario está dispuesto a volver una y otra vez. La tercera es el pago: el usuario está dispuesto a asignarle un presupuesto real. Completar solo la primera etapa todavía no constituye un negocio estable.
\end{importantbox}

\begin{warningbox}{No confundir difusión con ingresos}
La difusión atrae atención, pero los ingresos provienen de resolver repetidamente problemas reales. Si una aplicación de IA solo busca que la compartan, es fácil que el siguiente demo más vistoso la eclipse; solo si entra en el flujo de trabajo tiene la oportunidad de convertirse en una partida presupuestaria.
\end{warningbox}

\begin{knowledgebox}{Tabla de indicadores de PMF}
\begin{tabularx}{\textwidth}{p{0.22\textwidth}p{0.34\textwidth}X}
\toprule
Indicador & Descripción & Significado para Lovart \\
\midrule
Activación & El usuario produce por primera vez una obra que lo satisface & Demuestra que el Magic Moment existe. \\
Retención & Si el usuario vuelve una y otra vez para hacer proyectos nuevos & Demuestra que no es un juguete de un solo uso. \\
Profundidad & Si el usuario importa su marca, sus materiales y su colaboración & Demuestra que entró en un flujo de trabajo real. \\
Pago & Si las personas o los equipos están dispuestos a pagar de forma continua & Demuestra que su valor tiene precio. \\
Difusión & Si el usuario muestra sus obras y lo recomienda por iniciativa propia & Demuestra que la emoción y el producto se amplifican entre sí. \\
\bottomrule
\end{tabularx}
\end{knowledgebox}

\begin{warningbox}{El diagnóstico erróneo del PMF emocional}
Los productos de IA generan con facilidad una retroalimentación emocional intensa, pero esa retroalimentación puede deberse a la novedad y no a un valor de largo plazo. Al evaluar el PMF, es indispensable desglosar el «asombro del usuario» en retención, tasa de finalización de tareas, conversión a pago y profundidad de la colaboración en equipo.
\end{warningbox}

\begin{knowledgebox}{De usuarios gratuitos a usuarios profesionales}
Las aplicaciones de IA suelen atraer primero a muchos usuarios curiosos, pero el valor comercial de un Agent vertical proviene de los usuarios profesionales. Los usuarios profesionales exigen estabilidad, control, colaboración y entregables, y también están más dispuestos a pagar. Lo esencial para Lovart no es que todos lo prueben una vez, sino que los diseñadores lo incorporen al flujo de sus proyectos.
\end{knowledgebox}

\subsection{Resumen de la sección}

«Qué placer» es la percepción del PMF, pero no puede quedarse en una percepción. Las aplicaciones de IA deben convertir el pico emocional en un flujo de trabajo estable y en un ciclo comercial cerrado.
